\documentclass[11pt,a4paper]{article}
\usepackage{templates/modeling-notes}
\hypersetup{pdftitle={Task sheet 1 -- Complementarity and two-phase simplex}}
\title{\textbf{Mathematical Modeling}\\[2mm]
\large Task Sheet 1: Complementarity and Two-Phase Simplex}
\author{Personal study notes}
\date{}
\begin{document}
\maketitle
\tableofcontents

\section{How to read the certificates and tableaus}\label{sec:notation}
Feasibility comes before optimality. A proposed point must satisfy every row
and every variable-domain restriction. Complementary products alone cannot
certify an inadmissible point.

For a maximization problem, multiply each $\leq$ row by a nonnegative
weight and each $\geq$ row by a nonpositive weight. Summing produces
$y^TAx\leq y^Tb$. With $x\geq0$, the column conditions $A^Ty\geq c$
then give $c^Tx\leq b^Ty$: the dual minimizes a valid upper bound.
For minimization with $Ax\geq b$, $x\geq0$, nonnegative weights and
$A^Ty\leq c$ instead give the lower bound $b^Ty\leq c^Tx$.
If a variable is nonpositive, its column comparison reverses when multiplied
by that variable. The signs follow from these comparisons, rather than from
a rule to reverse every symbol indiscriminately.

Once primal and dual points are feasible, equal objective values prove
optimality. Complementarity helps find the dual point: a strict primal row
forces its weight to zero, and a nonzero primal coordinate forces its dual
column inequality to be an equality. For mixed rows, use their nonnegative
slacks and the corresponding signed weights consistently. These ideas are
developed in
\href{../lectures/04-degeneracy-and-duality.pdf}{Degeneracy, duality, and dual
  simplex}, in the sections on weighted bounds and complementary slackness.

For each simplex calculation below, use a maximization problem $\max d^Tv$
with $Mv=b$, $v\geq0$. An ordered basis $B$ produces
\[
  \bar M=M_B^{-1}M,\qquad \bar b=M_B^{-1}b,
  \qquad q_j=d_B^T\bar m_j-d_j,\qquad z_0=d_B^T\bar b.
\]
Tableau rows mean $\bar Mv=\bar b$ and $z+q^Tv=z_0$. The $c_B$ column gives the objective coefficient of each basic
variable. The $z_j$ row sums each constraint column with these weights;
subtract the displayed $c_j$ row to obtain $q_j=z_j-c_j$. The same
weighted sum on the RHS gives $z_0$. In Phase I these are the auxiliary
objective coefficients; in Phase II they are the original ones.
A negative $q_k$ offers improvement; increasing $v_k$ by $\theta$ changes
the basic coordinates to $\bar b-\theta\bar m_k$.
Only positive entries of $\bar m_k$ limit the step, with
$\theta^*=\min_{i:\bar m_{ik}>0}\bar b_i/\bar m_{ik}$.
Choose a most negative $q_k$, breaking entering ties by the smaller column
index. The displayed paths have no leaving tie.
The derivations appear in
\href{../lectures/03-simplex-method.pdf}{The simplex method}, in the sections
on canonical equations, one simplex step, and initialization.

\section{Task 2: complementary slackness}\label{sec:task-two}
\begin{otherlanguage}{polish}
  Korzystając z zasady komplementarności sprawdzić czy podane rozwiązanie jest
  rozwiązaniem optymalnym. Podać rozwiązanie dualne.
\end{otherlanguage}

\subsection{Task 2(a): a feasible candidate and a matching upper bound}
\paragraph{Task statement.}
\[
  \max(4x_1+5x_2-8x_3),\qquad
  \begin{cases}
    4x_1+x_2-2x_3\leq1, \\
    2x_1+x_2-2x_3\leq6, \\
    -x_1+x_2-x_3\leq2,  \\
    x_1+x_2-x_3\leq3,   \\
    x_1,x_2,x_3\geq0,
  \end{cases}
  \qquad x=(0,3,1).
\]
The candidate is nonnegative, and its row values are $(1,1,2,2)^T$.
The slacks are $(0,5,0,1)^T$, so it is feasible, with value $7$.
There is one nonnegative dual weight per upper-bound row:
\[
  \begin{aligned}
    \min\quad&y_1+6y_2+2y_3+3y_4,\\
    4y_1+2y_2-y_3+y_4&\geq4,\\
    y_1+y_2+y_3+y_4&\geq5,\\
    -2y_1-2y_2-y_3-y_4&\geq-8,\qquad y\geq0.
  \end{aligned}
\]
Strict primal rows force $y_2=y_4=0$. Since $x_2,x_3>0$, the last
two dual columns must be tight:
\[
  y_1+y_3=5,\qquad -2y_1-y_3=-8.
\]
Thus $y=(3,0,2,0)^T$. The remaining column equals $10\geq4$;
all dual signs hold. Its value is $3+2\cdot2=7$, matching the primal.
The dual surplus is $(6,0,0)^T$, so all complementary products vanish.
Both points are optimal.

This also proves uniqueness of the primal optimum. The positive first-column
surplus forces $x_1=0$; the positive weights $y_1,y_3$ force rows 1 and 3
to be tight. Solving $x_2-2x_3=1$, $x_2-x_3=2$ gives $(x_2,x_3)=(3,1)$.

\subsection{Task 2(b): reject the supplied point, then solve the dual}
\paragraph{Task statement.}
\[
  \min(6x_1+2x_2-3x_3),\qquad
  \begin{cases}
    x_1-x_2-x_3\geq-1, \\
    -3x_2-2x_3\geq-4,  \\
    x_1-x_2+x_3\geq4,  \\
    x_1,x_2,x_3\geq0,
  \end{cases}
  \qquad x=(2,0,-2).
\]
The supplied point has $x_3=-2<0$ and its third row is $0<4$.
It is infeasible, hence not optimal. We retain the stated point; replacing
its sign would answer a different candidate question.

The task also requests a dual solution. Nonnegative weights on the three
lower-bound rows yield
\[
  \begin{aligned}
    \max\quad&-y_1-4y_2+4y_3,\\
    y_1+y_3&\leq6,\\
    -y_1-3y_2-y_3&\leq2,\\
    -y_1-2y_2+y_3&\leq-3,\qquad y\geq0.
  \end{aligned}
\]
To find a feasible primal witness, the second row gives
$3x_2+2x_3\leq4$, hence $x_3\leq2$. The third row gives
$x_1\geq4+x_2-x_3$, and therefore
\[
  6x_1+2x_2-3x_3\geq24+8x_2-9x_3\geq6.
\]
The point $\hat x=(2,0,2)^T$ attains $6$; its row values are $(0,-4,4)^T$,
so it is feasible. Its positive first-row surplus forces $y_1=0$.
Its positive coordinates $\hat x_1,\hat x_3$ force
\[
  y_3=6,\qquad -2y_2+y_3=-3,
\]
giving $y=(0,9/2,6)^T$. Its second column is $-39/2\leq2$;
the other columns are tight. The dual value is $-18+24=6$.
Thus this is an optimal dual solution, and $\hat x$ is the actual primal
optimum. Equality in the elementary bound above forces $x_2=0$, $x_3=2$,
and $x_1=2$, proving primal uniqueness.

\subsection{Task 2(c): mixed row signs and an optimal segment}
\paragraph{Task statement.}
\[
  \max(2x_1+x_2-3x_3+x_4),\qquad
  \begin{cases}
    -x_1+x_2+x_3+x_4\geq3,    \\
    2x_1+x_2+2x_3+2x_4\leq10, \\
    x_1+x_2-x_3+x_4\leq2,     \\
    x_1,x_2,x_3,x_4\geq0,
  \end{cases}
  \qquad x=(0,5/2,1/2,0).
\]
The row values are $(3,7/2,2)^T$. The nonnegative slacks, in the
order ``left minus right'' for row 1 and ``right minus left'' for rows 2--3,
are $(0,13/2,0)^T$. The candidate is feasible and has value $1$.

For an upper bound, row 1 needs a nonpositive weight, whereas rows 2--3
need nonnegative weights:
\[
  \begin{aligned}
    \min\quad&3y_1+10y_2+2y_3,\\
    -y_1+2y_2+y_3&\geq2,\\
    y_1+y_2+y_3&\geq1,\\
    y_1+2y_2-y_3&\geq-3,\\
    y_1+2y_2+y_3&\geq1,\\
    y_1&\leq0,\qquad y_2,y_3\geq0.
  \end{aligned}
\]
The strict second row forces $y_2=0$. Positive $x_2,x_3$ require
$y_1+y_3=1$ and $y_1-y_3=-3$, hence $y=(-1,0,2)^T$.
The four dual column values are $(3,1,-3,1)^T$, so the surplus is
$(1,0,0,0)^T$. Its value is $-3+4=1$: both points are optimal.

The positive first-column surplus forces every primal optimum to have
$x_1=0$. Nonzero $y_1,y_3$ force rows 1 and 3 to be tight, giving
$x_3=1/2$ and $x_2+x_4=5/2$. The second row then equals $7/2+x_4$,
which remains below $10$ for $0\leq x_4\leq5/2$. Thus the entire optimal set is
\[
  x(t)=(0,5/2-t,1/2,t)^T,\qquad 0\leq t\leq5/2.
\]
The supplied point is one endpoint, not the unique optimum. A zero dual
surplus alone would not prove this; the displayed nonzero feasible changes do.

\section{Task 3: dual problems and two-phase simplex}\label{sec:task-three}
\begin{otherlanguage}{polish}
  Podaj ZD$(y)$. Rozwiąż metodą 2-fazową.
\end{otherlanguage}
Here ZD denotes the dual problem. Artificial variables provide a feasible
basis for an auxiliary system; they are not allowed to change the final model.
Phase I maximizes minus their sum, which is never positive. An optimum of
zero allows their removal and a return to the original objective. A negative
optimum would prove the original constraints infeasible.

\subsection{Task 3(a): a surplus variable enters during Phase I}
\paragraph{Task statement.}
\[
  \max(-x_1-2x_2+x_3),\qquad
  \begin{cases}
    x_1+x_3\leq12,     \\
    x_1+3x_2+x_3\geq4, \\
    x_1+2x_2-x_3\geq6, \\
    x_1,x_2,x_3\geq0.
  \end{cases}
\]
\paragraph{Dual from an upper bound.}
The first row receives $y_1\geq0$ and the two lower-bound rows receive
$y_2,y_3\leq0$. The dual is
\[
  \begin{aligned}
    \min\quad&12y_1+4y_2+6y_3,\\
    y_1+y_2+y_3&\geq-1,\\
    3y_2+2y_3&\geq-2,\\
    y_1+y_2-y_3&\geq1,\\
    y_1&\geq0,\qquad y_2,y_3\leq0.
  \end{aligned}
\]

\paragraph{Initialization.}
Add slack $s_1$, surpluses $s_2,s_3$, and artificials $a_2,a_3$:
\[
  \begin{aligned}
    x_1+x_3+s_1&=12,\\
    x_1+3x_2+x_3-s_2+a_2&=4,\\
    x_1+2x_2-x_3-s_3+a_3&=6.
  \end{aligned}
\]
All eight variables are nonnegative. Use column order
$(x_1,x_2,x_3,s_1,s_2,s_3,a_2,a_3)$, initial basis $B=(4,7,8)$,
and Phase I objective $w=-a_2-a_3$.
Surplus columns point in the negative unit direction and cannot supply
nonnegative initial basic values for positive right-hand sides.

\paragraph{Phase I: initial tableau.}
Substitution gives $w=-10+2x_1+5x_2-s_2-s_3$; the indicator row must
already account for both basic artificial variables.
\[
  \setlength{\arraycolsep}{3pt}
  \begin{array}{r@{\enspace}c@{\quad}rrrrrrrr@{\qquad}r@{\qquad\quad}l}\cmidrule(r{1.5em}){1-11}
    c_B & B & x_1 & x_2 & x_3 & s_1 & s_2 & s_3 & a_2 & a_3 & \bar b & \text{Ratio }(x_2\uparrow)\\
    & c_j & 0 & 0 & 0 & 0 & 0 & 0 & -1 & -1 &  & \\
    \cmidrule(r{1.5em}){1-11}
    0 & s_1 & 1 & 0 & 1 & 1 & 0 & 0 & 0 & 0 & 12 & \text{--}\\
    -1 & a_2 & 1 & 3 & 1 & 0 & -1 & 0 & 1 & 0 & 4 & 4 / 3\\
    -1 & a_3 & 1 & 2 & -1 & 0 & 0 & -1 & 0 & 1 & 6 & 6 / 2=3\\
    \cmidrule(r{1.5em}){1-11}
    & z_j & -2 & -5 & 0 & 0 & 1 & 1 & -1 & -1 & -10 & \\
    & q & -2 & -5 & 0 & 0 & 1 & 1 & 0 & 0 & -10 & \\
    \cmidrule(r{1.5em}){1-11}
  \end{array}
\]
Enter $x_2$. Ratios are $4/3$ and $6/2=3$, so $a_2$ leaves with
$\theta=4/3$. Divide its row by $3$ and eliminate column $x_2$:
\[
  \setlength{\arraycolsep}{3pt}
  \begin{array}{r@{\enspace}c@{\quad}rrrrrrrr@{\qquad}r@{\qquad\quad}l}\cmidrule(r{1.5em}){1-11}
    c_B & B & x_1 & x_2 & x_3 & s_1 & s_2 & s_3 & a_2 & a_3 & \bar b & \text{Ratio }(s_2\uparrow)\\
    & c_j & 0 & 0 & 0 & 0 & 0 & 0 & -1 & -1 &  & \\
    \cmidrule(r{1.5em}){1-11}
    0 & s_1 & 1 & 0 & 1 & 1 & 0 & 0 & 0 & 0 & 12 & \text{--}\\
    0 & x_2 & 1/3 & 1 & 1/3 & 0 & -1/3 & 0 & 1/3 & 0 & 4/3 & \text{--}\\
    -1 & a_3 & 1/3 & 0 & -5/3 & 0 & 2/3 & -1 & -2/3 & 1 & 10/3 & (10 / 3) / (2 / 3)=5\\
    \cmidrule(r{1.5em}){1-11}
    & z_j & -\tfrac{1}{3} & 0 & \tfrac{5}{3} & 0 & -\tfrac{2}{3} & 1 & \tfrac{2}{3} & -1 & -\tfrac{10}{3} & \\
    & q & -1/3 & 0 & 5/3 & 0 & -2/3 & 1 & 5/3 & 0 & -10/3 & \\
    \cmidrule(r{1.5em}){1-11}
  \end{array}
\]
Now $s_2$ has the most negative indicator. Its only positive entry is $2/3$
in row $a_3$, so $a_3$ leaves with $\theta=(10/3)/(2/3)=5$.
There is no rule restricting entry to original decision variables: a surplus
is a legitimate variable in the equivalent system.
\[
  \setlength{\arraycolsep}{3pt}
  \begin{array}{r@{\enspace}c@{\quad}rrrrrrrr@{\qquad}r}
    c_B & B & x_1 & x_2 & x_3 & s_1 & s_2 & s_3 & a_2 & a_3 & \bar b\\
    & c_j & 0 & 0 & 0 & 0 & 0 & 0 & -1 & -1 & \\
    \hline
    0 & s_1 & 1 & 0 & 1 & 1 & 0 & 0 & 0 & 0 & 12\\
    0 & x_2 & 1/2 & 1 & -1/2 & 0 & 0 & -1/2 & 0 & 1/2 & 3\\
    0 & s_2 & 1/2 & 0 & -5/2 & 0 & 1 & -3/2 & -1 & 3/2 & 5\\
    \hline
    & z_j & 0 & 0 & 0 & 0 & 0 & 0 & 0 & 0 & 0\\
    & q & 0 & 0 & 0 & 0 & 0 & 0 & 1 & 1 & 0\\
    \hline
  \end{array}
\]
Phase I has value zero and both artificials are nonbasic. Delete their
columns; the remaining basis $B=(4,2,5)$ already belongs to the original system.

\paragraph{Phase II: restore the original objective.}
Recompute the indicator row for the original coefficients:
\[
  d=(-1,-2,1,0,0,0)^T.
\]
\[
  \begin{array}{r@{\enspace}c@{\quad}rrrrrr@{\qquad}r}
    c_B & B & x_1 & x_2 & x_3 & s_1 & s_2 & s_3 & \bar b\\
    & c_j & -1 & -2 & 1 & 0 & 0 & 0 & \\
    \hline
    0 & s_1 & 1 & 0 & 1 & 1 & 0 & 0 & 12\\
    -2 & x_2 & 1/2 & 1 & -1/2 & 0 & 0 & -1/2 & 3\\
    0 & s_2 & 1/2 & 0 & -5/2 & 0 & 1 & -3/2 & 5\\
    \hline
    & z_j & -1 & -2 & 1 & 0 & 0 & 1 & -6\\
    & q & 0 & 0 & 0 & 0 & 0 & 1 & -6\\
    \hline
  \end{array}
\]
All indicators are nonnegative, so no Phase II pivot is needed.
The primal solution is $x^*=(0,3,0)^T$, with slacks/surpluses $(12,5,0)^T$
and maximum $-6$. Indeed, the original third constraint immediately gives
$-x_1-2x_2+x_3\leq-6$.
The dual point $y^*=(0,0,-1)^T$ has the required signs, satisfies every
column with equality, and has value $-6$.

There are alternative optima: $x(t)=(t,3-t/2,0)^T$ is feasible and has
value $-6$ for every $0\leq t\leq6$. This can also be read from the zero
$x_1$ indicator and its positive allowable step from the final basis.

\subsection{Task 3(b): a nonpositive variable and negative right-hand sides}
\paragraph{Task statement.}
\[
  \min(2x_1-2x_2+x_3),\qquad
  \begin{cases}
    x_1-x_2+x_3\geq-1,   \\
    x_1+2x_2+2x_3\geq-4, \\
    x_1-x_2-x_3\geq4,    \\
    x_1\geq0,\quad x_2\geq0,\quad x_3\leq0.
  \end{cases}
\]
\paragraph{Dual in the original variables.}
All lower-bound rows have nonnegative weights. The first two variables are
nonnegative, so their column comparisons are $\leq$; the nonpositive $x_3$
requires $\geq$. Explicitly,
\[
  \begin{aligned}
    \max\quad&-y_1-4y_2+4y_3,\\
    y_1+y_2+y_3&\leq2,\\
    -y_1+2y_2-y_3&\leq-2,\\
    y_1+2y_2-y_3&\geq1,\qquad y\geq0.
  \end{aligned}
\]
The last comparison ensures that
$(y_1+2y_2-y_3)x_3\leq x_3$ because $x_3\leq0$.

\paragraph{Transform and initialize.}
Set $t=-x_3\geq0$, so the original objective becomes
$F=2x_1-2x_2-t$. We maximize its negative $z=-F=-2x_1+2x_2+t$.
Multiply the first two constraint rows by $-1$ to obtain positive
right-hand sides; their inequality directions reverse:
\[
  \begin{aligned}
    -x_1+x_2+t+s_1&=1,\\
    -x_1-2x_2+2t+s_2&=4,\\
    x_1-x_2+t-s_3+a_3&=4.
  \end{aligned}
\]
Here $s_1,s_2$ are slacks of the reversed rows, $s_3$ a surplus, and
$a_3$ artificial. Use column order $(x_1,x_2,t,s_1,s_2,s_3,a_3)$ and
initial basis $B=(4,5,7)$. Phase I maximizes $w=-a_3$.
The two negative original right-hand sides do not require artificials after
the row reversal; their slack columns provide a feasible start.

\paragraph{Phase I.}
The initial tableau is
\[
  \begin{array}{r@{\enspace}c@{\quad}rrrrrrr@{\qquad}r@{\qquad\quad}l}\cmidrule(r{1.5em}){1-10}
    c_B & B & x_1 & x_2 & t & s_1 & s_2 & s_3 & a_3 & \bar b & \text{Ratio }(x_1\uparrow)\\
    & c_j & 0 & 0 & 0 & 0 & 0 & 0 & -1 &  & \\
    \cmidrule(r{1.5em}){1-10}
    0 & s_1 & -1 & 1 & 1 & 1 & 0 & 0 & 0 & 1 & \text{--}\\
    0 & s_2 & -1 & -2 & 2 & 0 & 1 & 0 & 0 & 4 & \text{--}\\
    -1 & a_3 & 1 & -1 & 1 & 0 & 0 & -1 & 1 & 4 & 4 / 1=4\\
    \cmidrule(r{1.5em}){1-10}
    & z_j & -1 & 1 & -1 & 0 & 0 & 1 & -1 & -4 & \\
    & q & -1 & 1 & -1 & 0 & 0 & 1 & 0 & -4 & \\
    \cmidrule(r{1.5em}){1-10}
  \end{array}
\]
Both $x_1$ and $t$ have indicator $-1$; the smaller column index chooses
$x_1$. Its only positive coefficient is in row $a_3$, so $a_3$ leaves
with $\theta=4$:
\[
  \begin{array}{r@{\enspace}c@{\quad}rrrrrrr@{\qquad}r}
    c_B & B & x_1 & x_2 & t & s_1 & s_2 & s_3 & a_3 & \bar b\\
    & c_j & 0 & 0 & 0 & 0 & 0 & 0 & -1 & \\
    \hline
    0 & s_1 & 0 & 0 & 2 & 1 & 0 & -1 & 1 & 5\\
    0 & s_2 & 0 & -3 & 3 & 0 & 1 & -1 & 1 & 8\\
    0 & x_1 & 1 & -1 & 1 & 0 & 0 & -1 & 1 & 4\\
    \hline
    & z_j & 0 & 0 & 0 & 0 & 0 & 0 & 0 & 0\\
    & q & 0 & 0 & 0 & 0 & 0 & 0 & 1 & 0\\
    \hline
  \end{array}
\]
The zero Phase I value and nonbasic artificial allow its column to be removed.

\paragraph{Phase II.}
For the negative original objective use $d=(-2,2,1,0,0,0)^T$:
\[
  \begin{array}{r@{\enspace}c@{\quad}rrrrrr@{\qquad}r@{\qquad\quad}l}\cmidrule(r{1.5em}){1-9}
    c_B & B & x_1 & x_2 & t & s_1 & s_2 & s_3 & \bar b & \text{Ratio }(t\uparrow)\\
    & c_j & -2 & 2 & 1 & 0 & 0 & 0 &  & \\
    \cmidrule(r{1.5em}){1-9}
    0 & s_1 & 0 & 0 & 2 & 1 & 0 & -1 & 5 & 5 / 2\\
    0 & s_2 & 0 & -3 & 3 & 0 & 1 & -1 & 8 & 8 / 3\\
    -2 & x_1 & 1 & -1 & 1 & 0 & 0 & -1 & 4 & 4 / 1=4\\
    \cmidrule(r{1.5em}){1-9}
    & z_j & -2 & 2 & -2 & 0 & 0 & 2 & -8 & \\
    & q & 0 & 0 & -3 & 0 & 0 & 2 & -8 & \\
    \cmidrule(r{1.5em}){1-9}
  \end{array}
\]
Enter $t$. The ratios are $5/2$, $8/3$, and $4$, so $s_1$ leaves
with $\theta=5/2$. Normalize its row and eliminate the $t$ column:
\[
  \begin{array}{r@{\enspace}c@{\quad}rrrrrr@{\qquad}r}
    c_B & B & x_1 & x_2 & t & s_1 & s_2 & s_3 & \bar b\\
    & c_j & -2 & 2 & 1 & 0 & 0 & 0 & \\
    \hline
    1 & t & 0 & 0 & 1 & 1/2 & 0 & -1/2 & 5/2\\
    0 & s_2 & 0 & -3 & 0 & -3/2 & 1 & 1/2 & 1/2\\
    -2 & x_1 & 1 & -1 & 0 & -1/2 & 0 & -1/2 & 3/2\\
    \hline
    & z_j & -2 & 2 & 1 & \tfrac{3}{2} & 0 & \tfrac{1}{2} & -\tfrac{1}{2}\\
    & q & 0 & 0 & 0 & 3/2 & 0 & 1/2 & -1/2\\
    \hline
  \end{array}
\]
All indicators are nonnegative, proving the maximum of $z$ is $-1/2$.
Recover $x_3=-t$ and $F=-z$: the original solution and minimum are
\[
  x^*=(3/2,0,-5/2)^T,\qquad F^*=1/2.
\]
Its original row values are $(-1,-7/2,4)^T$, with surpluses $(0,1/2,0)^T$.
The positive second surplus forces $y_2=0$. Nonzero $x_1,x_3$ force
$y_1+y_3=2$ and $y_1-y_3=1$, yielding $y^*=(3/2,0,1/2)^T$.
It is dual feasible, including the tight second column, and has value
$-3/2+2=1/2$.

The final objective row also gives
$F=1/2+(3/2)s_1+(1/2)s_3$.
Optimality therefore forces $s_1=s_3=0$. The remaining equations allow
\[
  x(\tau)=(3/2+\tau,\tau,-5/2)^T,\qquad \tau\geq0,
\]
with $s_2=1/2+3\tau$. This is the entire optimal set, an unbounded ray
with a finite constant objective. An unbounded feasible set does not imply
an unbounded objective.

\section{Checks and questions for understanding}\label{sec:checks}
All results above are exact rational calculations. The companion MATLAB
script checks the original variable domains, row residuals, primal--dual
objective equality, complementary products, and every displayed tableau.
It follows the specified bases and verifies entering choices, ratio tests,
step lengths, and the original-coordinate recovery; it is a check of these
hand solutions rather than a general LP solver.
Its absolute tolerance is $10^{-10}$ for these small, modestly scaled inputs.
No units are specified by the task statements; objective values are reported
without invented physical units.

\begin{enumerate}[leftmargin=*]
  \item Which two independent feasibility tests reject the supplied point in 2(b)?
  \item Why is $y_1\leq0$ in 2(c), but $y_1\geq0$ in 3(b)?
  \item Why may a surplus enter the basis in 3(a)?
  \item Why must the objective row be rebuilt between phases?
  \item Explain separately the row reversal, variable substitution, and
        objective negation in 3(b). Which changes the final reported objective sign?
  \item Use the displayed segment or ray to prove alternatives exist;
        explain why a zero indicator alone is insufficient.
\end{enumerate}

\section{Sources and editorial remarks}\label{sec:sources}
\textbf{Task source:} \texttt{ProgMatc1.pdf}, titled
\foreignlanguage{polish}{Programowanie Matematyczne zestaw zadań 1},
PDF page 1, group 2, and PDF page 2, group 3, supplied in the
\texttt{execrices-2026-27/} materials directory. The Polish instructions,
coefficients, inequality directions, variable domains, and supplied candidates
are retained in Sections~\ref{sec:task-two}--\ref{sec:task-three}.
English explanations, dual derivations, complete tableaus, and descriptions
of the optimal sets are added study material. No lecture date is inferred.

\paragraph{Numerical verification.}
The
\href{run:matlab/exercises/sheet_01_complementarity_and_two_phase.m}
{companion MATLAB script} can be run from the repository root with
\begin{lstlisting}
run('matlab/exercises/sheet_01_complementarity_and_two_phase.m')
\end{lstlisting}
All assertions passed in MATLAB R2026b, version 26.2.0.3386108.
The user supplied the output on 2026-10-11 and confirmed this version.
All complementary certificates and displayed tableaus passed at absolute
tolerance $10^{-10}$. The tableau calculations and certificates were also
independently checked with exact rational arithmetic.

Execution evidence: \path{docs/matlab-verification.md}.
\end{document}
